\documentclass[a4paper]{article}
\usepackage[a4paper, left=3cm, right=2cm, top=2cm, bottom=2cm]{geometry}
\usepackage{ragged2e}
\usepackage{fontspec}
\setmainfont{Times New Roman}
\usepackage{float}
\usepackage{setspace}
\singlespacing
\setlength{\parskip}{3pt}
\setlength{\parindent}{0pt}
\usepackage[fontsize=13pt]{fontsize}
\usepackage{url}

\usepackage{tikz}
\usetikzlibrary{
	arrows.meta,
	calc,
	shapes,
	shapes.misc,
	shapes.arrows,
	chains,
	matrix,
	positioning,
	scopes,
	decorations.pathmorphing,
	shadows
}

\usepackage{cite}
\usepackage{amsmath}
\usepackage{graphicx}
\usepackage{booktabs}
\usepackage{array}
\usepackage{url}
\usepackage{listings}
\graphicspath{{./}{paper/}{paper/figures/}}
\lstdefinestyle{mystyle}{
    backgroundcolor=\color{gray!10},   
    commentstyle=\color{green!60!black},
    keywordstyle=\color{blue},
    numberstyle=\tiny\color{gray},
    stringstyle=\color{purple},
    basicstyle=\ttfamily\footnotesize,
    breaklines=true,                   % Automatically wrap long lines
    captionpos=b,                      % Caption at the bottom
    keepspaces=true,                   % Keeps spaces for indentation
    numbers=left,                      % Line numbers on the left
    numbersep=5pt,                     
    showspaces=false,                  
    showstringspaces=false,
    showtabs=false,                    
    tabsize=2
}
\lstset{style=mystyle}

\usepackage{titlesec}
\usepackage{tocloft}
\usepackage[labelsep=period]{caption}

% Section numbering: PART I, PART II, ... and subsections I.1, I.2, ...
\renewcommand{\thesection}{\Roman{section}}
\renewcommand{\thesubsection}{\thesection.\arabic{subsection}}
\titleformat{\section}{\normalfont\Large\bfseries}{PART \thesection:}{0.5em}{}
\titleformat{\subsection}{\normalfont\normalsize\bfseries}{\thesubsection.}{0.5em}{}

% Caption names
\renewcommand{\figurename}{Figure}
\renewcommand{\tablename}{Table}

% Headings of the front-matter lists
\renewcommand{\contentsname}{TABLE OF CONTENTS}
\renewcommand{\listfigurename}{LIST OF FIGURES}
\renewcommand{\listtablename}{LIST OF TABLES}

% Centered, bold list titles
\renewcommand{\cfttoctitlefont}{\hfill\normalsize\bfseries}
\renewcommand{\cftaftertoctitle}{\hfill\null}
\renewcommand{\cftloftitlefont}{\hfill\normalsize\bfseries}
\renewcommand{\cftafterloftitle}{\hfill\null}
\renewcommand{\cftlottitlefont}{\hfill\normalsize\bfseries}
\renewcommand{\cftafterlottitle}{\hfill\null}
\setlength{\cftbeforetoctitleskip}{0pt}
\setlength{\cftaftertoctitleskip}{6pt}
\setlength{\cftbeforeloftitleskip}{18pt}
\setlength{\cftafterloftitleskip}{6pt}
\setlength{\cftbeforelottitleskip}{18pt}
\setlength{\cftafterlottitleskip}{6pt}

% TOC entries: "PART I: ...", bold, dotted leaders
\renewcommand{\cftsecfont}{\bfseries}
\renewcommand{\cftsecpagefont}{\bfseries}
\renewcommand{\cftsecpresnum}{PART~}
\renewcommand{\cftsecaftersnum}{:}
\renewcommand{\cftsecleader}{\bfseries\cftdotfill{\cftdotsep}}
\setlength{\cftsecnumwidth}{6.5em}
\setlength{\cftbeforesecskip}{4pt}

\renewcommand{\cftsubsecfont}{\bfseries}
\renewcommand{\cftsubsecpagefont}{\normalfont}
\renewcommand{\cftsubsecaftersnum}{.}
\setlength{\cftsubsecindent}{1em}
\setlength{\cftsubsecnumwidth}{5.5em}

% List of figures: "Figure 1. ..."
\renewcommand{\cftfigpresnum}{Figure~}
\renewcommand{\cftfigaftersnum}{.}
\setlength{\cftfigindent}{0pt}
\setlength{\cftfignumwidth}{4.8em}

% List of tables: "Table 1. ..."
\renewcommand{\cfttabpresnum}{Table~}
\renewcommand{\cfttabaftersnum}{.}
\setlength{\cfttabindent}{0pt}
\setlength{\cfttabnumwidth}{4.8em}

\begin{document}
	
	\title{Báo cáo kết quả nghiên cứu \\
		{\large Memo-Lang: Ngữ nghĩa đọc hủy trên các khối lượng công việc ma trận GPU}}
	
	\author{
		Dự án trình thông dịch độc lập\\
		Tháng 9 năm 2026, Hà Nội}
	
	% ================= PAGE 1: COVER PAGE =================
	\thispagestyle{plain}
	\begin{center}
		\bfseries
		CUỘC THI KHOA HỌC KỸ THUẬT CẤP TRƯỜNG DÀNH CHO HỌC SINH TRUNG HỌC, NĂM HỌC 2026 -- 2027
	\end{center}
	
	\vspace{5cm}
	
	\begin{center}
		\bfseries
		BÁO CÁO KẾT QUẢ NGHIÊN CỨU
	\end{center}
	
	\vspace{3cm}
	
	\begin{center}
		{\Large\bfseries\itshape Tên dự án:}\\[10pt]
		{\Large\bfseries
		\MakeUppercase{MEMOLANG: NGHIÊN CỨU VÀ PHÁT TRIỂN CƠ CHẾ GIẢI PHÓNG VÙNG ĐỆM CHO PHÁT TRIỂN AI}\par}
	\end{center}
	
	\vspace{3cm}
	
	\noindent\textbf{\textit{Lĩnh vực:} Phần mềm hệ thống}
	
	\vspace{6pt}
	
	\vfill
	
	\begin{center}
		\textbf{\textit{Hà Nội, tháng 9 năm 2026}}
	\end{center}
	
	\newpage
	
	% ================= PAGES 2–3: TOC + LISTS =================
	\tableofcontents
	
	\listoffigures
	\listoftables
	
	% ----- LIST OF ABBREVIATIONS -----
    \vspace{10cm}
	\vspace{18pt}
	\begin{center}
		\textbf{DANH MỤC TỪ VIẾT TẮT}
	\end{center}
	\vspace{-6pt}
	\begin{center}
		\renewcommand{\arraystretch}{1.15}
		\begin{tabular}{>{\bfseries}l >{\raggedright\arraybackslash}p{5.8cm} >{\raggedright\arraybackslash}p{5.8cm}}
			\toprule
			\textbf{Từ viết tắt} & \textbf{Tên đầy đủ} & \textbf{Mô tả} \\
			\midrule
			AI      & Artificial Intelligence & Hệ thống máy tính thực hiện các tác vụ thường đòi hỏi trí tuệ con người \\
			AOT     & Ahead-Of-Time (compilation) & Biên dịch thành mã máy trước khi chương trình chạy \\
			AST     & Abstract Syntax Tree & Cây biểu diễn cấu trúc của mã nguồn \\
			cuBLAS  & CUDA Basic Linear Algebra Subprograms & Thư viện đại số tuyến tính tăng tốc bằng GPU của NVIDIA \\
			CUDA    & Compute Unified Device Architecture & Nền tảng tính toán song song của NVIDIA cho GPU \\
			CUTLASS & CUDA Templates for Linear Algebra Subroutines & Thư viện mẫu C++ của NVIDIA cho đại số tuyến tính trên GPU \\
			DRAM    & Dynamic Random-Access Memory & Bộ nhớ chính cần được làm tươi định kỳ \\
			FFI     & Foreign Function Interface & Cơ chế gọi hàm được viết bằng ngôn ngữ khác \\
			FP32    & 32-bit Floating Point & Định dạng số dấu phẩy động độ chính xác đơn \\
			GEMM    & General Matrix Multiplication & Quy trình chuẩn để nhân hai ma trận \\
			GPU     & Graphics Processing Unit & Bộ xử lý chuyên dụng cho tính toán song song cao \\
			IR      & Intermediate Representation & Dạng mã trung gian trong trình biên dịch, nằm giữa mã nguồn và mã máy \\
			JIT     & Just-In-Time (compilation) & Biên dịch thành mã máy trong khi chương trình đang chạy \\
			LLM     & Large Language Model & Mô hình AI được huấn luyện trên lượng lớn văn bản \\
			LLVM    & LLVM Compiler Infrastructure (originally ``Low Level Virtual Machine'') & Bộ công cụ để xây dựng trình biên dịch và bộ sinh mã \\
			MatMul  & Matrix Multiplication & Phép nhân hai ma trận \\
			REPL    & Read--Eval--Print Loop & Vỏ tương tác thực thi từng dòng lệnh được nhập vào \\
			\bottomrule
		\end{tabular}
	\end{center}
	
	
	\newpage
	\maketitle
	
	\begin{abstract}
		\normalsize
		Sự phát triển của AI trong thời gian gần đây đã dẫn đến sự bùng nổ của nhiều thuật toán máy học và mô hình học sâu. Do vậy, nhu 
		cầu cho năng lực tính toán để học trên những bộ dữ liệu khổng lồ
		cũng từ đó tăng.
		
		Hiện nay, phép toán được sử dụng rộng rãi cho thuật toán máy học là GEMM (General Matrix Multiplication - Phép nhân ma trận tổng 
		quát). Trong quá trình học, mỗi phép nhân ma trận sẽ tạo ra một ma trận khác và ma trận mới sẽ được nối với những phép nhân ma trận khác. Để xử lý nhiều thao tác, tác vụ và quá trình khác nhau, những mô hình máy học cần thực hiện hàng triệu phép tính ma trận. Do những phép tính ma trận này so với các phép tính thông thường khác có tính đặc thù về thời gian rất cao, vì vậy những phép tính tiếp theo dần bị dồn lên bộ nhớ đệm và các bước trung gian cho đến khi những chường trình quản lý bộ nhớ hoặc quy tắc của máy tính giải phóng bộ nhớ.
		
		Việc lặp lại quá trình này đối với mô hình máy học trở thành một vấn đề lớn trong việc hoạt động dưới một khối lượng công việc rất lớn, để lại hậu quá lớn cho hiệu suất và độ trễ của mô hình trên thiết bị và tài nguyên phần cứng nó đang sử dụng.
		
		Bởi vậy, phần mềm đồng thời góp phần tối ưu hóa nhân của GEMM và tăng tuổi thọ của phần cứng được sử dụng, đảm bảo bộ đệm, dữ liệu đệm và con trỏ được giải phóng sớm nhất có thể và giảm phí tổn từ kết quả của việc dồn bộ nhớ đệm.
	\end{abstract}
	
	
	\section{Câu hỏi nghiên cứu}
	
	Mục tiêu chính của chúng tôi là liên kết trực tiếp \emph{thư viện chạy cuBLAS} đã biên dịch sẵn cho một ứng dụng cụ thể, được đánh giá bằng một phương pháp cụ thể, dưới một mô hình thực thi tương tự \emph{hệ thống kiểu tuyến tính chặt chẽ}. Một số ngôn ngữ khác cũng theo đuổi ngữ nghĩa theo dõi việc tiêu thụ giá trị: hệ thống quyền sở hữu của Rust~\cite{rust} áp đặt nguyên tắc dùng một lần tương tự thông qua hệ thống kiểu affine, còn Futhark~\cite{futhark} là ngôn ngữ mảng hàm thuần túy dùng kiểu duy nhất (uniqueness types) để cập nhật tại chỗ một cách an toàn. Tuy nhiên, Futhark tự sinh kernel GPU từ mã nguồn thay vì tích hợp với các thư viện hiệu năng cao bên ngoài do bên thứ ba duy trì như cuBLAS, và đây được xem là một trong những trở ngại thực tế của ngôn ngữ này. Ngôn ngữ của chúng tôi thay vào đó tách việc phân tích tiêu thụ lúc biên dịch khỏi việc thực thi kernel.
	
	\begin{equation}
		T_1 = A B,\qquad
		T_2 = T_1 C,\qquad
		T_3 = T_2 D,
	\end{equation}
	
	Trong triển khai thực tế, nhiều vùng đệm trung gian tồn tại đồng thời. Tốt nhất là vùng lưu trữ của các vùng đệm này được xóa và tái sử dụng ngay khi phép toán tương ứng đã hoàn tất hoặc khi vòng đời của đối tượng kết thúc.
	
	Có thể kỳ vọng lợi ích về quản lý bộ nhớ từ một ngôn ngữ có ngữ nghĩa đọc hủy. Các phép toán số học như phép nhân ma trận có thể được giao cho kernel CUTLASS/cuBLAS đã tối ưu để vượt qua các hạn chế của Futhark~\cite{futhark}. Việc triển khai này được kỳ vọng không ảnh hưởng đến hiệu năng tính toán so với triển khai trên các ngôn ngữ khác. Tuy nhiên, cải thiện toàn trình được kỳ vọng xuất hiện nhờ các tối ưu hóa của ngôn ngữ, bao gồm:
	\begin{itemize} 
		\item Giảm cấp phát tạm thời
		\item Rút ngắn vòng đời vùng đệm
		\item Giảm áp lực bộ nhớ GPU
	\end{itemize}
	Về lý thuyết, hiệu quả của cách triển khai này tăng theo khối lượng công việc, vì các tensor trung gian hoặc các phép toán ma trận nhỏ không chịu ảnh hưởng đáng kể từ việc quản lý bộ nhớ so với các chuỗi phép toán lớn.
	
	Tuy nhiên, cần cân nhắc chi phí phát sinh vốn có, cũng như khả năng đánh giá quá cao chi phí bộ nhớ hoặc tính thực tiễn của việc tối ưu bộ nhớ. Do đó, dự án khảo sát các giả thuyết sau:
	
	\begin{itemize}
		\item Ngữ nghĩa đọc hủy có thể giảm mức sử dụng bộ nhớ đỉnh đối với các khối lượng công việc chứa các ma trận trung gian tồn tại ngắn.
		\item Đọc hủy có thể giảm số lần hoặc chi phí cấp phát vùng đệm tạm khi các vùng đệm đã tiêu thụ có thể được tái sử dụng một cách an toàn.
		\item Với các khối lượng công việc mà quản lý bộ nhớ chiếm phần đáng kể trong tổng thời gian thực thi, việc tái sử dụng vùng đệm có thể cải thiện hiệu năng toàn trình.
	\end{itemize}
	
	Các lợi ích thực tiễn tiềm năng gồm:
	
	\begin{itemize}
		\item Cải thiện hiệu năng của phép nhân ma trận, hỗ trợ huấn luyện và suy luận nhanh hơn cho các mô hình học máy.
		\item Tái sử dụng vùng đệm tensor trung gian hiệu quả hơn nhờ trình biên dịch.
		\item Giảm chi phí cấp phát vùng đệm và bộ nhớ đệm khi chạy.
	\end{itemize}
	
	\subsection{Thuật toán}
	
	Phép nhân ma trận được viết như sau
	
	\begin{equation}
		C_{i,j} = \sum_{k=0}^{N-1} A_{i,k}B_{k,j}
	\end{equation}
	
	Mỗi giá trị của ma trận kết quả $C$ được tạo ra bằng cách nhân một hàng của $A$ với một cột của $B$ rồi cộng các tích. Các triển khai hiệu quả phải quản lý việc di chuyển dữ liệu, tái sử dụng bộ nhớ và, trong sơ đồ đọc hủy, vòng đời của các vùng đệm toán hạng sau khi chúng bị tiêu thụ.
	
	Đọc hủy là thao tác đọc bộ nhớ vốn xóa giá trị vừa đọc: dữ liệu bị hủy cùng lúc với khi được truy cập. Nếu một giá trị cần được giữ lại, phải có thao tác ghi lại khôi phục nó ngay lập tức, hoặc một vùng đệm/bộ nhớ đệm bên ngoài phải giữ một bản sao.
	
	Các nghiên cứu liên quan về hệ thống kiểu tuyến tính có cách tiếp cận tương tự ở cấp độ ngôn ngữ, yêu cầu mỗi biến được dùng đúng một lần.
	
	Về mặt lịch sử, đọc hủy phổ biến trong thiết kế phần cứng và mạch điện hơn là trong phần mềm. Dybdahl và cộng sự~\cite{DybdahlKGN06} đã nghiên cứu hành vi đọc hủy trong DRAM nhúng và phát hiện rằng việc bỏ bộ đệm ghi lại, thay vào đó dựa vào bộ nhớ đệm (đọc dữ liệu vào bộ nhớ đệm mà không khôi phục ngay về DRAM, và chỉ ghi lại khi dòng bộ nhớ đệm bị thay thế), mang lại lợi ích đo được về điện năng và hiệu năng. 
	
	Dù kết quả này riêng cho hệ thống bộ nhớ phần cứng, nó gợi ra câu hỏi rộng hơn mà dự án nghiên cứu: liệu một nguyên tắc tiêu thụ-khi-đọc tương tự, áp dụng ở cấp phần mềm/trình thông dịch cho các khối lượng công việc nặng về GEMM, có thể mang lại kết quả tương tự hay không?
	
	
	\section{Thiết kế và phương pháp}
	
	Mục tiêu của ngôn ngữ là xác định liệu ngữ nghĩa đọc hủy ở cấp ngôn ngữ có thể mang lại lợi ích quản lý bộ nhớ đo được cho các chương trình GPU thiên về ma trận, đồng thời vẫn giữ được hiệu năng của các kernel nhân ma trận đã tối ưu hay không.
	
	Dự án sẽ được triển khai và đánh giá thông qua các mục tiêu sau:
	
	\begin{itemize}
		\item Định nghĩa cú pháp lõi giống Python và ngữ nghĩa đọc hủy của ngôn ngữ.
		\item Xây dựng bộ phân tích từ vựng, bộ phân tích cú pháp, AST và quy trình sinh mã dựa trên LLVM.
		\item Sinh LLVM IR cho các phép toán số học và ma trận.
		\item Tích hợp phép nhân ma trận GPU đã tối ưu thông qua CuBLAS.
		\item Xây dựng cơ chế theo dõi quyền sở hữu và sự tiêu thụ tường minh cho các toán hạng ma trận.
		\item Xác định khi nào các vùng đệm ma trận đã tiêu thụ có thể được tái sử dụng một cách an toàn.
		\item Kiểm chứng kết quả sinh ra với một tham chiếu số học độc lập.
		\item Đo mức sử dụng bộ nhớ GPU đỉnh và hành vi cấp phát vùng đệm tạm.
		\item Đo cả thời gian thực thi kernel GPU và thời gian chạy toàn chương trình.
		\item So sánh thực thi đọc hủy với một triển khai tương đương không đọc hủy.
		\item So sánh ngôn ngữ với các mốc cơ sở là thư viện GPU đã tối ưu khi phù hợp.
		\item Lưu giữ nhật ký benchmark, kết quả kiểm tra tính đúng đắn, đầu ra của trình biên dịch và các sản phẩm build để kết quả thực nghiệm có thể tái lập.
	\end{itemize}
	
	Đánh giá sẽ không mặc định rằng đọc hủy mang lại cải thiện hiệu năng. Việc giảm mức sử dụng bộ nhớ đỉnh mà không giảm thời gian chạy sẽ được coi là một kết quả có ý nghĩa; còn nếu quan sát thấy thời gian chạy được cải thiện, sẽ phân tích xem nó xuất phát từ việc giảm chi phí quản lý bộ nhớ hay từ thay đổi ở chính kernel nhân ma trận.
	
	\subsection{Ngữ nghĩa đọc hủy}
	
	Đọc hủy coi một giá trị là đã bị tiêu thụ khi một phép toán có quyền sử dụng độc quyền giá trị đó. Sau khi bị tiêu thụ, giá trị gốc không còn khả dụng cho các phép toán tiếp theo, trừ khi lập trình viên chủ động tạo bản sao hoặc bảo toàn nó bằng cách khác.
	
	Với các phép toán ma trận, điều này tạo cơ hội tái sử dụng vùng lưu trữ. Ví dụ,
	
	\begin{lstlisting}
		C = A @ B
		D = C @ W
	\end{lstlisting}
	
	có thể cho phép vùng lưu trữ gắn với $A$ hoặc $C$ được tái sử dụng khi trình biên dịch xác định rằng giá trị tương ứng không còn nơi tiêu thụ nào nữa.
	
	Mục đích của cơ chế này không phải là thay đổi phép tính toán học mà phép nhân ma trận thực hiện. Thay vào đó, nó thay đổi vòng đời của các vùng đệm gắn với phép tính. Nếu một vùng đệm có thể được tái sử dụng mà không cần sao chép nội dung, việc triển khai có thể giảm mức tiêu thụ bộ nhớ đỉnh và hoạt động cấp phát.
	
	Việc tối ưu này vốn phụ thuộc vào khối lượng công việc. Một ma trận còn nhiều nơi tiêu thụ trong tương lai thì không thể bị hủy một cách an toàn sau lần dùng đầu tiên. Trong trường hợp đó, việc triển khai có thể cần sao chép, mượn hoặc một cơ chế khác để bảo toàn giá trị. Do đó, ngữ nghĩa đọc hủy có thể tạo ra sự đánh đổi giữa hiệu quả bộ nhớ và tính linh hoạt khi tái sử dụng giá trị.
	
	Vì vậy, ngôn ngữ coi đọc hủy là một cơ chế sở hữu tường minh, thay vì giả định rằng việc tiêu thụ mọi toán hạng chắc chắn cải thiện hiệu năng.
	
	\subsection{Kiến trúc trình biên dịch và thư viện chạy}
	
	\subsection{ Thiết kế cú pháp }
	
	Chúng tôi chọn quy ước cú pháp giống Python cho ngôn ngữ lập trình:
	\begin{itemize}
		\item Chú thích dùng dấu thăng.
		\item Kiểu nguyên thủy: số nguyên, số thực, luận lý và chuỗi.
		\item Khai báo biến.
		\item Hàm dùng 'def'.
		\item extern cho liên kết runtime/FFI (giữa các ngôn ngữ).
		\item Độ ưu tiên toán tử.
		\item Cấu trúc điều khiển.
	\end{itemize}
	Tuy nhiên, vì tính khả thi khi triển khai và các quyết định thiết kế tiếp theo, chúng tôi điều chỉnh một số điểm của cú pháp Python gốc:
	\begin{itemize}
		\item độ ưu tiên đặc biệt của clone() để sao chép bộ nhớ
		\item dùng dấu ngoặc nhọn thay cho thụt lề để xác định phạm vi hàm
	\end{itemize}
	
	\subsection{Sơ đồ Free-on-Consume và thư viện chạy cuBLAS}
	Sơ đồ được nêu trong \cite{DybdahlKGN06} sẽ được sử dụng với triết lý thiết kế đã điều chỉnh để hoạt động ở cấp phần mềm.
	\begin{table}[htbp]
  \centering
  \caption{Các lựa chọn triển khai sơ đồ trong thiết kế.}
  \label{tab:scheme_implementation}
  \small
  \begin{tabular}{@{}>{\RaggedRight}p{0.48\linewidth} >{\RaggedRight}p{0.48\linewidth}@{}}
    \toprule
    \textbf{Sơ đồ} & \textbf{Triển khai} \\
    \midrule
    Dữ liệu trước hết được đọc từ DRAM vào bộ nhớ đệm. 
      & Không khác biệt. \\ \addlinespace
    
    Dữ liệu được ghi lại vào DRAM khi dòng bộ nhớ đệm bị thay thế, bất kể dữ liệu có bị sửa đổi hay không. 
      & Các toán hạng được giải phóng vô điều kiện, với giả định bộ nhớ đệm luôn bị sửa đổi. \\ \addlinespace
    
    Tái sử dụng kiến trúc bộ nhớ đệm hiện có. 
      & Dùng cấp phát theo thoát phạm vi để quản lý bộ nhớ đệm. \\ \addlinespace
    
    Che giấu độ trễ ghi lại bằng chồng lấp nhiều bank (xen kẽ bank). 
      & Che giấu chi phí giải phóng bộ nhớ sau các lần khởi chạy kernel GPU. \\
    \bottomrule
  \end{tabular}
\end{table}

	Về việc liên kết cuBLAS đã biên dịch sẵn:
	\begin{itemize}
		\item cuBLAS nên được liên kết động.
		\item Trình điều khiển nên liên kết và đóng gói cùng các tệp nhị phân đã biên dịch của cuBLAS.
		\item Để thuận tiện, được phân phối dưới dạng một tệp đóng gói.
	\end{itemize}

	
	\subsection{ Đường ống xử lý } 
	\begin{figure}[H]
	\centering
	\begin{tikzpicture}[
		node distance=1.1cm and 1.3cm,   % vertical and horizontal gap between box borders
		box/.style={
			rectangle, draw, rounded corners,
			text width=6.6cm,
			minimum height=2.8cm,
			inner sep=5pt,
			align=center,
			font=\normalsize
		},
		special/.style={box, fill=gray!20},
		arr/.style={-{Latex[length=2mm]}, thick}
		]
		
		\node[box] (src) {\textbf{Mã nguồn}\\
			Văn bản thô viết theo cú pháp của ngôn ngữ; lỗi cú pháp được chuyển đến bộ chẩn đoán.};
		\node[box, right=of src] (lex) {\textbf{Bộ phân tích từ vựng}\\
			Tạo luồng token gồm các bộ (loại, giá trị, vị trí).};
		\node[box, below=of lex] (par) {\textbf{Bộ phân tích cú pháp}\\
			Phân tích đệ quy xuống kết hợp độ ưu tiên toán tử để xây dựng AST và báo lỗi.};
		\node[box, left=of par] (ast) {\textbf{AST}\\
			Cây các nút biểu thức và câu lệnh biểu diễn cấu trúc chương trình.};
		\node[special, below=of ast] (con) {\textbf{Phân tích tiêu thụ [Đọc hủy]}\\
			Đặt \texttt{consume\_a}/\texttt{consume\_b}; xem phần sau.};
		\node[box, right=of con] (ir) {\textbf{LLVM IR}\\
			Sinh các lời gọi thư viện chạy kèm cờ tiêu thụ.};
		\node[box, below=of ir] (jit) {\textbf{JIT / AOT}\\
			Biên dịch IR thành mã máy bằng bộ công cụ LLVM.};
		\node[box, left=of jit] (link) {\textbf{Liên kết}\\
			Liên kết mã đã biên dịch với thư viện chạy nhân ma trận đã biên dịch sẵn thành tệp thực thi có thể chạy trên GPU.};
		\node[special, below=of link] (gpu) {\textbf{Kernel GPU [Đọc hủy]}\\
			Thực hiện nhân ma trận và giải phóng các vùng đệm đã tiêu thụ.};
		
		\draw[arr] (src) -- (lex);
		\draw[arr] (lex) -- (par);
		\draw[arr] (par) -- (ast);
		\draw[arr] (ast) -- (con);
		\draw[arr] (con) -- (ir);
		\draw[arr] (ir) -- (jit);
		\draw[arr] (jit) -- (link);
		\draw[arr] (link) -- (gpu);
	\end{tikzpicture}
	\caption{Đường ống biên dịch của ngôn ngữ.}
	\label{fig:pipeline}
\end{figure}
			

	Phần front-end của trình biên dịch được viết bằng C++; mỗi giai đoạn được thiết kế trước, sau đó lập trình, rồi gỡ lỗi độc lập để tìm lỗi, cuối cùng mới lắp ráp theo từng mô-đun.
	\section{Quá trình, xây dựng và kiểm thử}
	\subsection{Quá trình}

		\begin{table}[h]
		\centering
		\caption{Các tệp mã nguồn chính}
		\label{tab:modules}
		\begin{tabular}{ll}
			\toprule
			Tệp & Mục đích \\
			\midrule
			\texttt{Lexer.h/.cpp} & Tách token mã nguồn của ngôn ngữ \\
			\texttt{Parser.h/.cpp} & Xây dựng AST \\
			\texttt{ExprAST.h} & Định nghĩa các nút AST \\
			\texttt{CodeGen.cpp} & Sinh LLVM IR \\
			\texttt{CuBLASRuntime.cpp} & Liên kết (binding) MatMul của cuBLAS \\
			\texttt{main.cpp} & Trình điều khiển / REPL \\
			\bottomrule
		\end{tabular}
	\end{table}
	
	
	
	
	\subsection{Kiểm thử}
	
	\begin{table}[h]
		\centering
		\caption{Kết quả kiểm thử tính đúng đắn}
		\label{tab:verification}
		\begin{tabular}{ll}
			\toprule
			Trường hợp kiểm thử & Kết quả \\
			\midrule
			MatMul $N \times N$ so với tham chiếu (thông thường) & Chưa có \\
			Phát hiện lỗi dùng lại giá trị đã đọc hủy (khi chạy) & Chưa có \\
			Trường hợp biên (rỗng, 1x1, không vuông) & Chưa có \\
			\bottomrule
		\end{tabular}
	\end{table}

	MatMul $N \times $N so với tham chiếu (thông thường): Đánh giá một chuỗi phép nhân ma trận để kiểm tra chịu tải nhằm đánh giá cải thiện toàn trình.
	Phát hiện lỗi dùng lại giá trị đã đọc hủy: Đánh giá khi chạy và khi biên dịch để phân loại lỗi.
	Trường hợp biên: Đánh giá các trường hợp biên để kiểm tra chịu tải về tính tương thích và độ toàn diện của thuật toán.
	
	
	\subsection{Đánh giá hiệu năng (Benchmark)}
	
	Việc đánh giá hiệu năng được thiết kế để phân biệt cải thiện về quản lý bộ nhớ với cải thiện của chính phép nhân ma trận.
	
	Ba triển khai chính sẽ được so sánh:
	
	\begin{enumerate}
		\item Một triển khai thư viện GPU đã tối ưu, dùng làm tham chiếu hiệu năng kernel.
		\item Ngôn ngữ dùng ngữ nghĩa giá trị thông thường (không đọc hủy).
		\item Ngôn ngữ dùng ngữ nghĩa đọc hủy.
	\end{enumerate}
	
	Khi phù hợp, một triển khai C++ hoặc CUDA đơn giản cũng sẽ được đưa vào làm mốc cơ sở.
	
	Phép so sánh quan trọng nhất là giữa hai cấu hình của ngôn ngữ. Cả hai cấu hình thực hiện cùng các phép toán toán học và dùng cùng một triển khai nhân ma trận GPU bên dưới. Phép so sánh này cô lập ảnh hưởng của ngữ nghĩa đọc hủy tốt hơn so với việc so sánh với một triển khai không liên quan.
	
	Bộ benchmark sẽ gồm các phép nhân ma trận đơn lẻ cũng như các chuỗi phép toán ma trận. Kích thước ma trận sẽ gồm nhỏ, vừa, lớn và không vuông. Điều này cho phép đánh giá xem ảnh hưởng của đọc hủy có phụ thuộc vào quy mô khối lượng công việc hay không.
	
	Các phép đo sau sẽ được thu thập:
	
	\begin{itemize}
		\item Thời gian thực thi kernel GPU.
		\item Thời gian thực thi toàn trình.
		\item Mức tiêu thụ bộ nhớ GPU đỉnh.
		\item Số lần cấp phát vùng đệm tạm.
		\item Số lần tái sử dụng vùng đệm.
		\item Thời gian truyền từ máy chủ sang thiết bị và ngược lại, khi có áp dụng.
		\item Chi phí phát sinh của trình biên dịch và thư viện chạy.
		\item Sai số số học so với triển khai tham chiếu.
		\item Công suất tiêu thụ của GPU.
	\end{itemize}
	
	Các phép đo GPU sẽ dùng nhiều lần lặp và các vòng khởi động phù hợp. Chi phí khởi tạo CUDA \cite{CUDA} và biên dịch sẽ được tách khỏi các phép đo thực thi ở trạng thái ổn định khi có thể.
	
	Benchmark sẽ báo cáo cả số đo tuyệt đối và mức thay đổi tương đối. Không có cải thiện hiệu năng nào được quy cho ngữ nghĩa đọc hủy trừ khi các số đo tương ứng ủng hộ cách diễn giải đó.
	
	
	
	\subsection{Mối quan hệ kỳ vọng giữa ngôn ngữ và hiệu năng MatMul}
	
	Thời gian thực thi kỳ vọng có thể được phân tách gần đúng như sau
	
	\begin{equation}
		T_{\mathrm{total}} =
		T_{\mathrm{language}} +
		T_{\mathrm{memory}} +
		T_{\mathrm{launch}} +
		T_{\mathrm{kernel}}.
	\end{equation}
	
	Ngữ nghĩa đọc hủy chủ yếu nhắm vào $T_{\mathrm{memory}}$. Chúng không được kỳ vọng làm giảm đáng kể $T_{\mathrm{kernel}}$ khi dùng cùng một kernel GPU đã tối ưu.
	
	Với một phép nhân ma trận lớn, $T_{\mathrm{kernel}}$ có thể chiếm phần lớn tổng thời gian thực thi, dẫn đến cải thiện toàn trình rất ít quan sát được. Tuy nhiên, với các khối lượng công việc chứa nhiều tensor trung gian hoặc các phép toán ma trận nhỏ, $T_{\mathrm{memory}}$ và các chi phí chạy khác có thể trở nên đáng kể hơn. 

	Việc đánh giá hiệu quả của ngôn ngữ chủ yếu dựa vào việc đo $T_{\mathrm{memory}}$.
	
	
	
	
	\subsection{Bằng chứng và khả năng tái lập}
	\label{sec:evidence}
	
	Ví dụ biên dịch: 
	\begin{lstlisting}[caption = Ví dụ biên dịch]
	def basicadd(int i, int k) {
		return i + k;
		}
	# Basic compilation of source code, addition only. No showcase of GEMMs or cuBLAS integration.
	\end{lstlisting}
	\section{Kết luận}

	\subsection{Tiến độ}
	Dự án đã hoàn thành một phần công việc đã lên kế hoạch, bao gồm bộ phân tích từ vựng, bộ phân tích cú pháp, AST và một trình điều khiển cơ bản. Thiết kế cho phần còn lại của đường ống đã được chốt, nhưng việc triển khai các tính năng ở phía sau của đường ống chưa bắt đầu do hạn chế về thời gian.

\subsection{Công việc tương lai}
	Các tính năng ở phía sau của đường ống, bao gồm việc tích hợp thư viện chạy nhân ma trận đã biên dịch sẵn, vẫn cần được triển khai. Các tính năng đang được lên kế hoạch như sau:
	\begin{itemize}
		\item Kiểu vector và ma trận
		\item Trình điều khiển chẩn đoán và bộ máy chẩn đoán
		\item Phân tích tiêu thụ tĩnh (kiểm tra đọc hủy) và xử lý các vùng đệm và giá trị trung gian của kernel GPU
		\item Liên kết với thư viện chạy nhân ma trận đã biên dịch sẵn để tạo tệp thực thi có khả năng khởi chạy kernel GPU
		\item Tách phần sinh mã khỏi AST để việc hạ cấp mã sạch hơn
		\item Các bước tối ưu hóa
		\item Hệ thống macro (khai báo liên kết, trừu tượng hóa mẫu chuỗi nhân ma trận)
\end{itemize}

	\subsection{Ý nghĩa}
		Dự án này vừa là một thí nghiệm nhằm khám phá các cải tiến tiềm năng cho phép toán GEMM trong các thuật toán học máy và AI, bổ sung cho thư viện cuBLAS \cite{cublas} vốn đã được tối ưu cao, vừa là một phản hồi trước một vấn đề rộng hơn của lĩnh vực này. Việc huấn luyện mô hình hiện nay tiêu tốn lượng bộ nhớ khổng lồ, và nhu cầu từ hạ tầng AI đã vượt quá nguồn cung chip nhớ toàn cầu, góp phần khiến giá DRAM và NAND tăng kéo dài, gây áp lực tài chính lên các công ty CNTT, các tổ chức nghiên cứu khoa học và người tiêu dùng \cite{DRAM_PRICE}. 
		 Các thuật toán tính toán GEMM đã được tối ưu đến \emph{gần giới hạn lý thuyết} về thông lượng thô; nghiên cứu này thay vào đó nhắm vào hiệu quả bộ nhớ, dùng một hệ thống sở hữu mới để giảm dung lượng bộ nhớ cần thiết cho mỗi khối lượng công việc, như một đóng góp nhằm giảm bớt áp lực từ phía cầu vốn đang gây ra tình trạng thiếu hụt bộ nhớ hiện nay.

	\subsection{Tính mới}
		\begin{enumerate}
			\item Sản phẩm được xây dựng dựa trên những ý tưởng đã có như kiểu tuyến tính và quyền sở hữu. Tính mới của dự án nằm ở việc áp dụng các ý tưởng đó để giải phóng và quản lý vùng đệm GPU, bằng cơ chế đọc hủy được kiểm tra tĩnh khi biên dịch. 
			\item Dự án tiếp cận đề tài theo hướng vừa thử nghiệm với \emph{hệ thống kiểu tuyến tính chặt chẽ} vừa để chế tạo.
		\end{enumerate}

		

	\begin{thebibliography}{9}
		\bibitem{DybdahlKGN06}
		Haakon Dybdahl, Per Gunnar Kjeldsberg, Marius Grann{\ae}s, and Lasse Natvig,
		``Destructive-read in embedded DRAM, impact on power consumption,''
		\emph{Journal of Embedded Computing}, vol.~2, no.~1, pp.~83--96, 2006.
		
		\bibitem{cublas}
		NVIDIA, ``cuBLAS Library.'' \url{https://developer.nvidia.com/cublas}
		
		\bibitem{llvm}
		LLVM Project, ``The LLVM Compiler Infrastructure.'' \url{https://llvm.org}

		\bibitem{rust}
		N.~D.~Matsakis and F.~S.~Klock, ``The Rust language,'' in \emph{Proc. ACM SIGAda Ada Letters}, 2014. \url{https://www.rust-lang.org}
		
		\bibitem{futhark}
		T.~Henriksen, N.~G.~W.~Serup, M.~Elsman, F.~Henglein, and C.~E.~Oancea, ``Futhark: purely functional GPU-programming with nested parallelism and in-place array updates,'' in \emph{Proc. 38th ACM SIGPLAN Conference on Programming Language Design and Implementation (PLDI)}, 2017. \url{https://futhark-lang.org}

		\bibitem{Mixed-precision training}
		NVIDIA, ''Mixed-precision training.'' \url{https://docs.nvidia.com/deeplearning/performance/mixed-precision-training/index.html}

		\bibitem{nvidia_accuracy}
		NVIDIA, ``Accuracy Considerations''
		\url{https://docs.nvidia.com/deeplearning/tensorrt/latest/inference-library/accuracy-considerations.html},

		\bibitem{DRAM_PRICE}
		TrendForce, ``Memory Price Outlook for 1Q26 Sharply Upgraded,''
		\url{https://www.trendforce.com/presscenter/news/20260202-12911.html},

		\bibitem{CUDA}
		NVIDIA, ``CUDA ''
		\url{https://developer.nvidia.com/cuda/toolkit}

	\end{thebibliography}
	
\end{document}